\section{Datasets}
The datasets contained within the GPAT modelling framework help maintain a consistent flow of data between models, providing information about the specific flights in the run and the associated fuel burn and emissions masses throughout those flights (FL). The dispersion and advection of their corresponding plumes is captured in the Plume Dataset (PL). To investigate the ensuing climate effects of those emissions through gridded chemical and microphysical analysis, the Background Chemistry (BG\_CHEM) and Meteorology (MET) Datasets is defined. The Emissions Dataset (EMI) is generated through aggregation of 1D emissions and plume geo-vector data to a 2D Eulerian grid. The Contrail (CONTRAIL) and Chemistry (CHEM) Datasets are then generated by concatenating background chemistry, meteorology and emissions datasets, followed by chemical and microphysical processing through the run\_cc and run\_boxm models. This final step provides insight into how the emissions present in the model affect the composition of the surrounding atmosphere, over the specified runtime of the simulation. All data structures contained within the GPAT framework are created using the Pandas, Xarray and pycontrails Python libraries~\cite{Xarray2024, pycontrails2024}. 

\subsection{Flight Dataset (FL)}
The Flight Dataset is first generated when the user-defined Flight Parameters (fl\_params) are passed through the Trajectory Model. The latitude, longitude, altitude coordinates as well as the waypoint number are calculated for each time step throughout the flight, according to t0\_fl, ts\_fl and rt\_fl. Flight attributes are also assigned to this dataset as metadata; the flight's unique identifier (UID), aircraft type, number of engines, wingspan, maximum Mach number, maximum altitude and engine UID.\@

The Flight Dataset is passed to the Aircraft Performance Model, which estimates aircraft fuel flow using the aircraft positional data, aircraft-specific characteristics and the instantaneous meteorological situation. The down-selected meteorology data is assigned prior to the performance run. Aircraft performance is then simulated, and the dataset is updated with columns for true airspeed, thrust, thrust setting, rate of climb or descent, engine efficiency, fuel flow, fuel burn and aircraft mass at each flight waypoint. Total fuel burn is also added to the FL dataset's attributes. An example of a fully populated Flight Dataset is provided in table~\ref{FL_table}. 

% FL table

\subsection{Meteorology Dataset (MET)}
MET provides information on the meteorological situation over all time steps and grid cells of the simulation region, derived from the outputs of the GPAT Meteorological Model (gen\_met). The air temperature, water vapour concentration and wind field data are taken from the UK Met Office Unified Model (UM) for the year 1998 at a resolution of 5\textdegree~$\times$~5\textdegree~\cite{}. The pressure is estimated according to the pressure levels and associated altitude bands from the model output levels in Collins et al. 1997~\cite{Collins1997}, presented in table~\ref{Model_output_levels}. 

% Model output levels table

The Meteorological Model then calculates additional parameters, including specific humidity, relative humidity, molar mass of air (M) as well as atmospheric concentrations of oxygen (\ce{O2}) and nitrogen (\ce{N2}). The model also calculates the Solar Zenith Angle (SZA) across all grid cells and time steps in the simulation domain. Table~\ref{MET_table} shows an example Met Dataset. 

% Met table

\subsection{Background Chemistry Dataset (BG\_CHEM)}
This dataset contains the necessary chemical concentration data for all the background trace species initialised in the BOXM Photochemical Trajectory Model (see~\ref{run_boxm_section}). These data are obtained from the STOCHEM Common Representative Intermediates (CRI) v2.2 R5 model outputs presented in Khan et al.~\cite{Khan2020}. The Background Chemistry Model downloads the species data based on the simulation start time. The STOCHEM-CRI model outputs are at a 5\textdegree~$\times$~5\textdegree~resolution, and are monthly averaged, requiring further interpolation to fit the simulation domain. Table~\ref{BG_CHEM_table} shows all species and example values.

% BG CHEM table

\subsection{Plume Dataset (PL)}
The Plume Dataset is a container for all data associated with the plume waypoints from every flight in the simulation

\subsection{Emissions Dataset (EMI)}

\subsection{Contrail Dataset (CONTRAIL)}

\subsection{Chemistry Dataset (CHEM)}



